\SourceLeaf{82}
\noindent\hspace*{\parindent}Let $\mathfrak{P}$ denote the required function. We then have:

\[
\begin{aligned}\left(\frac{d\mathfrak{P}}{dP}\right) ={}& \text{constant part, or the part free of } q \text{, in }\\ &\frac{q\,.\,f'q\,.\,f\dfrac{P^{\frac1n}}{Q}}{nP\left(1-\left(\dfrac{P}{Q}\right)^{\frac{\nu}{n}}\right)}.
\end{aligned}
\]



\Needspace{28\baselineskip}\begin{center}
REMARKS ON THE NOTES ON SERIES INVERSION.
\end{center}

\noindent\hspace*{\parindent}The first note is taken from a letter to \textsc{Olbers} dated 12.\,Dec. 1813. The expansion coefficient denoted there by the symbol $(n,\nu)$ is given explicitly by the multinomial theorem in the form

\[
(n,\nu) = \sum_\lambda \frac{(n+\mu-1)(n+\mu-2).\,.\,.(n+1)n}{\lambda'!\,\lambda''!\,\lambda'''!\,.\,.\,.\,.}\, a^{\prime\lambda'}a^{\prime\prime\lambda''}a^{\prime\prime\prime\lambda'''}.\,.\,.,
\]

\noindent when $a'=-\alpha$, $a''=-\beta$, $a'''=-\gamma$, .\,.\,. are substituted. The sum extends over all sets of nonnegative integers $\lambda'$, $\lambda''$, .\,.\,., satisfying the equation $\lambda'+2\lambda''+3\lambda'''+.\,.\,.=\nu$, and $\mu=\lambda'+\lambda''+\lambda'''+.\,.\,.$.

The second note, which \textsc{Gauss} wrote in his copy of »\textsc{Hobert} und \textsc{Ideler}, trigonometrische Tafeln« (first page of the cover), gives explicitly the generalization indicated at the end of the first note. The inversion formula in question can easily be verified by means of the expression for $(n,\nu)$ just derived.

In the third note, which is on a separate slip and bears the date 1820 Jun.\,22, the generalization of the first note to arbitrary powers of $x$ is again developed explicitly by a slightly modified method. The inference from the two representations of $\Omega$ to equations I, II, .\,.\,. rests on a consideration of homogeneity.

The fourth derivation probably dates from 1816 and is found in a handbook. Here the symbol $[Fx]_{\text{Coeff.}\,x^n}$ means the coefficient of $x^n$ in the expansion of $F(x)$ in powers of $x$. \hfill\textsc{Fricke.}

